\begin{theorem}[Orthonormal coordinates for weighted bond coefficients]%
    \label{thm:peps_theta_bond_orthonormal_coordinates}
    \lean{TNLean.PEPS.exists_thetaBond_orthonormalCoordinates}
    \leanok
    \uses{thm:peps_representation_theta_squared, thm:peps_multiplicity_bond_weight}
    Let a representation $\rho$ of a finite group on a finite-dimensional
    complex Hilbert space have a
    supplied orthogonal internal decomposition into irreducible subrepresentations $S_i$,
    each with character multiplicity one. Choose orthonormal bases of the $S_i$, and let $D_i(g)$ be the matrices of
    the restricted representations in these coordinates.
    There is one global orthonormal basis $e$ for which
    \begin{align}
        [\rho(g)]_e&=\bigoplus_iD_i(g),&
        [\Theta^2\rho(g)]_e&=\bigoplus_i\sqrt{d_i}\,D_i(g).
    \end{align}
    Apply the normalized multiplicity-restoring bond map to the actual
    diagonal blocks of the second matrix. Its result is
    $\bigl(D_i(g)\otimes I_{d_i}\bigr)_i$.
    This is the coefficient transformation displayed in
    \cite[Section~7]{Schuch2010PEPS}.
    The assertion concerns the occurring irreducible sectors in the supplied
    decomposition; it does not require all irreducibles to occur.
    Its coordinate identification is unitary. Orthogonality and the
    decomposition are supplied hypotheses; this auxiliary assertion does not
    construct them or identify complete PEPS.
\end{theorem}

\begin{proof}\leanok
    Collect the sector orthonormal bases into a global orthonormal basis, simultaneously
    for the endomorphisms $\rho(g)$ and $\Theta^2\rho(g)$.
    Each preserves every irreducible summand. On the summand $S_i$, the
    fourth-root square acts by $\sqrt{d_i}$, so the second family has precisely
    the stated weighted diagonal blocks. Irreducibility gives $d_i>0$.
    The normalized bond map sends
    $\sqrt{d_i}D_i(g)$ to $D_i(g)\otimes I_{d_i}$, cancelling its normalization.
\end{proof}
